\scanpage{5}
\[
\begin{aligned}
\varphi(u_W(w)ar_0,i)&=u_W(\varphi(wr_0,i))\qquad\text{ou simplement}\\
\deleted{\unsure{\varphi(u(w)ar_0,i)=u(\varphi(wr_0,i))}}&\\
\Vert\qquad &\qquad\unsure{\leftarrow|\Vert}\\
\deleted{\unsure{\operatorname{int}(u(w))\,\varphi(ar_0,i)=u\varphi}} & u(w\,\varphi(r_0,i)\,w^{-1})\\
u(w)a\,\varphi(r_0,i)a^{-1}u(w)^{-1}&\qquad\Vert\\
&u(w)\,u(\varphi(r_0,i))\,u(w)^{-1}
\end{aligned}
\]
\deleted{donc} ou encore
\[
 a\,\varphi(r_0,i)a^{-1}=u\,\varphi(r_0,i)\qquad\forall\,i\in I
\]
i.e. $u_W=\operatorname{int}(a)$ (car les $\varphi(r_0,i)$, \inserted{$i\in I$,} engendrent $W$

donc
\[
u_R(wr_0)=(awa^{-1})(ar_0)=a(wr_0)=a_R(wr_0)
\]
i.e.
\[
u_R=a_R
\]
i.e.
\[
(u_W,u_R,\operatorname{id}_I)=\alpha(a).
\]

On a \deleted{\unsure{donc}} un diagramme de suites exactes
\diagram{\begin{tikzcd}[ampersand replacement=\&,column sep=small]
\& \& \& \substack{\text{\unsure{gr|gp|gr.} des autom ext}\\\text{de }\Sigma}\arrow[d,equal] \&\\
1\arrow[r] \& W\arrow[r,"\alpha_0"]\arrow[d,Rightarrow,"\mathrm{id}"] \& \operatorname{Aut}(\Sigma)\arrow[r,"\beta_0"]\arrow[d] \& \operatorname{Aut}(I)\arrow[r]\arrow[d,Rightarrow,"\mathrm{id}"] \& 1\\
1\arrow[r,dash] \& W\arrow[r,"\alpha"] \& \operatorname{Aut}(\Phi(\Sigma))\arrow[r,"\beta"] \& \operatorname{Aut}(I) \&
\end{tikzcd}}
et le diagramme des 5 \unsure{mor.} montre alors que la flèche verticale médiane est un iso, cqfd.

\begin{center}\rule{2.5cm}{0.35pt}\end{center}

Les \deleted{\unsure{parties}} \deleted{\unsure{ens}} paraboliques de racines de $\Sigma\subset V$ \inserted{(vectoriel réel de rang \unsure{$l$|1})} correspondent \deleted{\unsure{ent}} \deleted{aux} par \deleted{\unsure{orthogonalité}} \inserted{polarité} aux \uline{facettes} \inserted{\uline{fermées}} du syst. d'hyperplans dans $\check V$ associé au syst. de racines dans $V$.
% archive sheet number: 4

